\documentclass[11pt]{amsart}

\usepackage{amsmath}
\usepackage{amssymb}
\usepackage{amsthm}
\usepackage{url}
\usepackage{hyperref}
%\usepackage{psfig}

\begin{document}
\baselineskip=16pt
\textheight=8.5in
%\parindent=0pt 
\def\sk {\hskip .5cm}
\def\skv {\vskip .05cm}
\def\cos {\mbox{cos}}
\def\sin {\mbox{sin}}
\def\tan {\mbox{tan}}
\def\intl{\int\limits}
\def\lm{\lim\limits}
\newcommand{\frc}{\displaystyle\frac}
\def\xbf{{\mathbf x}}
\def\fbf{{\mathbf f}}
\def\gbf{{\mathbf g}}

\def\dbA{{\mathbb A}}
\def\dbB{{\mathbb B}}
\def\dbC{{\mathbb C}}
\def\dbD{{\mathbb D}}
\def\dbE{{\mathbb E}}
\def\dbF{{\mathbb F}}
\def\dbG{{\mathbb G}}
\def\dbH{{\mathbb H}}
\def\dbI{{\mathbb I}}
\def\dbJ{{\mathbb J}}
\def\dbK{{\mathbb K}}
\def\dbL{{\mathbb L}}
\def\dbM{{\mathbb M}}
\def\dbN{{\mathbb N}}
\def\dbO{{\mathbb O}}
\def\dbP{{\mathbb P}}
\def\dbQ{{\mathbb Q}}
\def\dbR{{\mathbb R}}
\def\dbS{{\mathbb S}}
\def\dbT{{\mathbb T}}
\def\dbU{{\mathbb U}}
\def\dbV{{\mathbb V}}
\def\dbW{{\mathbb W}}
\def\dbX{{\mathbb X}}
\def\dbY{{\mathbb Y}}
\def\dbZ{{\mathbb Z}}

\def\la{{\langle}}
\def\ra{{\rangle}}
\def\summ{{\sum\limits}}
\def\eps{{\varepsilon}}
\def\lam{{\lambda}}
\def\tr{{\rm tr}}
\def\char{{\rm char}}
\def\Mat{{\rm Mat}}
\def\Func{{\rm Func}}
\def\Span{{\rm Span}}
\def\Inv{{\rm Inv}}
\def\inv{{\rm inv}}
\def\Ker{{\rm Ker}}
\def\Im{{\rm Im}}
\def\rk{{\rm rk}}
\def\phi{{\varphi}}
\def\row{{\rm row}}
\def\col{{\rm col}}
\def\Rowspace{{\rm Rowspace}}
\def\null{{\rm null}}


\bf\centerline{Homework \#9. Due on Thursday, Nov 4th, 11:59pm on Canvas}\rm
\vskip .1cm

\bf{Reading for this homework: }\rm 5.1 and parts of 5.2

\bf{Plan for next week (Nov 2,4): }\rm  5.2 (Diagonalizability) continued,
7.3 (Minimal Polynomials) and 5.4 (Invariant subspaces and the Cayley-Hamilton Theorem) 

\skv
\noindent 

\bf\centerline{Problems: }\rm
\skv
{\bf Problem 1:} Let $F$ be a field.
Recall that the trace of a square matrix over $F$ denoted by $\tr$ is the sum of its diagonal entries.
Also recall from HW\#1.2 that $\tr(AB)=\tr(BA)$ for any $A,B\in \Mat_{n}(F)$. 
\begin{itemize}
\item[(a)] Prove that if $A,C\in \Mat_{n}(F)$ are similar (that is, $C=Q^{-1}AQ$ for some
invertible $Q\in \Mat_{n}(F)$), then $\tr(A)=\tr(C)$.
\item[(b)] Use (a) to define $\tr(T)$ for a linear transformation $T:V\to V$ where $\dim(V)<\infty$
(similarly to how we defined $\det(T)$ in Lecture~17). Your definition should not depend on the choice of a basis of $V$.
\end{itemize}
\skv
{\bf Problem 2:} Let $V$ be a vector space over a field $F$ and $T:V\to V$ a linear map. Given a polynomial $p(x)\in\mathcal{P}(F)$,
we can define $p(T)$ (the value of $p$ at $T$) in a natural way: if $p(x)=\sum\limits_{k=0}^n a_k x^k$ with $a_k\in F$,
we set $p(T)=\sum\limits_{k=0}^n a_k T^k$ (where $T^0=I$, the identity map). 
\begin{itemize}
\item[(a)] Prove that $p(T)$ is a linear map and $\Ker\, p(T)$ and $\Im\, p(T)$ are both $T$-invariant subspaces of $V$.
\item[(b)] Let $v\in V$ be an eigenvector for $T$ with eigenvalue $\lam$. Prove that $v$ is also an eigenvector for $p(T)$
with eigenvalue $p(\lam)$.
\end{itemize}
\skv
The results of Problems 3 and 4 below will be essential for the proof of the existence of the Jordan Canonical Forms of matrices
(tentatively scheduled for the Nov 9 week). 

\skv
{\bf Problem 3:} Let $V$ be a finite-dimensional vector space over a field $F$, $n=\dim(V)$ and $T:V\to V$ a linear map.
If $W$ is a $T$-invariant subspace of $V$, denote by $T_W: W\to W$ the map obtained from $T$ by restricting both the domain
and codomain to $W$ (we can restrict the codomain as well precisely because $W$ is $T$-invariant). 

Suppose now that $V=\bigoplus\limits_{i=1}^k W_i$ where each $V_i$ is $T$-invariant, and choose an ordered basis $\beta_i$ for each $V_i$.
Prove that 
\begin{itemize}
\item[(i)] $\beta=\beta_1\cup\ldots\cup\beta_k$ is a basis of $V$. In (ii) below we view $\beta$ as an ordered basis where elements of 
$\beta_1$ come first, followed by elements of $\beta_2$ etc.
\item[(ii)] The matrix $[T]_{\beta}$ is $k\times k$ block-diagonal, with $i^{\rm th}$ diagonal block equal to $[T_{W_i}]_{\beta_i}$,
that is,
$$[T]_{\beta}=\begin{pmatrix}
{[T_{W_1}]}_{\beta_1} & \ldots & 0\\
\vdots & \ddots & \vdots\\
0 & \ldots & {[T_{W_k}]}_{\beta_k}\\
\end{pmatrix}$$
\end{itemize}
{\bf Note:} Part (i) has nothing to do with $T$. Also (i) holds without the assumption $\dim(V)<\infty$, but in the finite-dimensional
case you can give a shorter proof using HW\#8.6.
\skv

\skv
{\bf Problem 4:} Let $V$ be a vector space and $T:V\to V$ a linear map. In parts (b)--(e) assume that $\dim\, V<\infty$.
\begin{itemize}
\item[(a)] Prove that $\Ker(T^n)\subseteq \Ker(T^{n+1})$ for each $n\geq 1$.
\item[(b)] By (a) we get an ascending chain of subspaces of $V$: $$\Ker(T)\subseteq \Ker(T^2)\subseteq \ldots .$$
Prove that this chain stabilizes after finitely many steps, that is, there is $N\in\dbN$ s.t. $\Ker(T^n)= \Ker(T^{N})$
for all $n\geq N$.
\item[(c)] Assume that $\Ker(T)\cap \Im(T)\neq \{0\}$. Prove that $\Ker(T^2)\neq \Ker(T)$.
\item[(d)] Use (b) and (c) (applied to a suitable power of $T$) to show that there exists $n\in\dbN$ s.t. $\Ker(T^n)\cap \Im(T^n)=\{0\}$. Deduce that $V=\Ker(T^n)\oplus \Im(T^n)$. 
\item[(e)] Use (d) to show that if $T$ is non-invertible and non-nilpotent, then $V$ can be written as a direct sum of 2 nonzero
$T$-invariant subspaces.
\end{itemize}
\skv
{\bf Problem 5:} Let $F$ be a field and $n\in\dbN$. Given an ordered basis $\beta$ of $F^n$, let $Q(\beta)$ be the matrix whose columns are the elements of $\beta$ (in the given order). Given $\lam_1,\ldots,\lam_n\in F$, let $D(\lam_1,\ldots,\lam_n)$ be the diagonal matrix with the diagonal entries $\lam_1,\ldots,\lam_n$ (in this order).

Prove that for a given $A\in \Mat_n(F)$, an ordered basis $\beta=\{v_1,\ldots, v_n\}$ of $F^n$ and scalars $\lam_1,\ldots,\lam_n\in F$,
the following are equivalent:
\begin{itemize}
\item[(i)] $v_i$ is an eigenvector of $A$ with eigenvalue $\lam_i$ for each $1\leq i\leq n$;
\item[(ii)] $QD=AQ$ (or, equivalently, $Q^{-1}AQ=D$) where $Q=Q(\beta)$ and $D=D(\lam_1,\ldots,\lam_n)$.
\end{itemize}
Deduce from this that $A$ is diagonalizable if and only if $A$ is similar to a diagonal matrix.
\skv
{\bf Hint:} Show that whenever the product of two matrices $XY$ is defined, the $j^{\rm th}$ column of $XY$ is equal to $X\col_j(Y)$ for all $j$.
Note that the result of Problem~5 is stated in the book as Corollary after Theorem~5.1 on page 247, but as far as I can tell, the proof is not clearly presented. 
\skv
{\bf Problem 6:} In each of the following examples the field is $\dbR$ (real numbers).
Find all eigenvalues and eigenvectors of the given $A\in \Mat_3(\dbR)$ and determine
whether the matrix $A$ is diagonalizable.
If $A$ is diagonalizable, find an invertible matrix $Q$
and a diagonal matrix $D$ s.t. $Q^{-1}AQ=D$.
\begin{align*}
&\mbox{(a)}\quad A=\begin{pmatrix} 1& 1& 0\\ 0& 1 &2 \\ 0& 0& 3\end{pmatrix}&
&\mbox{(b)}\quad A=\begin{pmatrix} 7& -4& 0\\8 &-5 &0\\ 6& -6& 3\end{pmatrix}&
\end{align*}
\skv
{\bf Problem 7:} Let $A=\Mat_n(F)$, and suppose that $\chi_T(x)$ splits completely over $F$,
that is, $\chi_T(x)=\prod\limits_{i=1}^n (x-\lam_i)$ for some $\lam_1,\ldots, \lam_n\in F$.
(Thus, $\lam_1,\ldots, \lam_n$ are precisely the eigenvalues of $T$ listed with their algebraic multiplicities).
Prove that $$\sum_{i=1}^n \lam_i=\tr(A)\quad\mbox{ and }\quad\prod_{i=1}^n \lam_i=\det(A).\eqno (***)$$
Then explain why the analogous results holds if $A$ is a replaced by a linear map $T:V\to V$ where
$V$ is a finite-dimensional vector space.
{\bf Hint:} Use HW\#8.5.
\skv
{\bf Problem 8:} Let $V$ be a vector space, and let $T:V\to V$ and $S:V\to V$ be linear maps which commute ($TS=ST$).
Let $\lam$ be an eigenvalue of $T$. Prove that the eigenspace $E_{\lam}(T)$ is $S$-invariant.
\end{document}

